\documentclass[a4paper,12pt]{article}

\usepackage[utf8]{inputenc}
\usepackage[T1]{fontenc}
\usepackage[table]{xcolor}
\usepackage{geometry}
\usepackage{longtable}
\usepackage{enumitem}
\usepackage{hyperref}
\usepackage{listings}
\usepackage{titlesec}
\usepackage{booktabs}
\usepackage{array}

\geometry{margin=2.5cm}

\definecolor{criticalcolor}{HTML}{DC3545}
\definecolor{highcolor}{HTML}{FD7E14}
\definecolor{mediumcolor}{HTML}{FFC107}
\definecolor{lowcolor}{HTML}{28A745}
\definecolor{resolvedcolor}{HTML}{6C757D}

\newcommand{\severitybox}[2]{%
  \colorbox{#1}{\makebox[2.5cm][c]{\textcolor{white}{\textbf{#2}}}}%
}

\newcommand{\fixedbox}{%
  \colorbox{resolvedcolor}{\makebox[2.5cm][c]{\textcolor{white}{\textbf{FIXED}}}}%
}

\lstset{
  basicstyle=\small\ttfamily,
  breaklines=true,
  frame=single,
  numbers=none,
  tabsize=2,
  backgroundcolor=\color{gray!10},
  extendedchars=true,
  inputencoding=utf8,
  literate={—}{---}1
}

\title{\textbf{Performance Report \#07}\\[4pt]
\large Vulkan Engine — Post-Refactoring Audit, Residual Bugs \& Performance}
\author{Henrique Rocha}
\date{\today}

\begin{document}
\maketitle

\begin{abstract}
Since report \#06, 19 commits have been merged, resolving 11 of 26 previously
identified issues (2 critical, 4 high, 5 medium). This report audits the residual
15 issues, adds 16 newly discovered findings (including latent race conditions,
API misuse, and resource management hazards), and provides an updated overall
assessment. Despite significant architectural improvements, several systemic
issues remain in the global state layer and validation handler.
\end{abstract}

\tableofcontents
\newpage

\section{Scope \& Methodology}

This audit covers the full codebase (\texttt{vulkan/}, \texttt{space/},
\texttt{services/}, \texttt{main.cpp}) with emphasis on the changes since
commit \texttt{c31bf96} (report \#06 baseline). Each issue was verified by
reading source code and cross-referencing with the Vulkan specification (1.3
\& 1.4) and the engine's own \texttt{AGENTS.md} guidelines.

\subsection*{Resolved Since Report \#06}

The following issues from perf\_report\_06 have been fixed:

\begin{center}
\begin{tabular}{lc}
\toprule
\textbf{ID} & \textbf{Fix Commit} \\
\midrule
B1 — Async ring wrap corruption & \texttt{e0315df} \\
B2 — Device-lost skips cleanup & \texttt{97da4b7} \\
B3 — Async upload buffer leak & \texttt{9c7be04} \\
B5 — Heap-allocated descriptor structs & \texttt{f9079e6} \\
P1 — Staging ring double-free & \texttt{77d4602} \\
P2 — Direct \texttt{vkAllocateMemory} & \texttt{4f945c7} \\
P4 — Map/unmap per operation & \texttt{64b98da} \\
S1 — Shadow cascade wrong stage mask & \textit{(commits around \texttt{c0d6e05})} \\
S2 — \texttt{ALL\_COMMANDS\_BIT} flush & \texttt{6a02101} \\
S3 — \texttt{deviceWaitIdle} drain & \texttt{0ae836b} \\
C1 — Legacy render pass parameter & \texttt{9c49551} \\
\midrule
\textbf{11 resolved} & \\
\bottomrule
\end{tabular}
\end{center}

The remaining 15 issues (6 open + 9 partially addressed) are carried forward
with updated status where applicable.

\newpage

\section{Issue Register}

\subsection{Critical Severity}

\begin{enumerate}[leftmargin=*]

\item[\textbf{[G1]}] \textbf{Global Mutable State Violates RAII Encapsulation}
  \hfill \severitybox{criticalcolor}{CRITICAL}

  \begin{tabular}{ll}
    \textbf{File:} & \texttt{vulkan/VulkanApp.cpp} \\
    \textbf{Line:} & 136--163 \\
    \textbf{Category:} & Architecture / Deadlock Risk \\
  \end{tabular}

  \textbf{Description:}
  Eight global data structures are declared at file scope outside the
  \texttt{VulkanApp} class:

  \begin{itemize}[noitemsep]
    \item \texttt{pendingCommandBuffers} (guarded by \texttt{pendingCmdMutex})
    \item \texttt{g\_cmdSubmitMap}
    \item \texttt{commandBufferPoolMap}
    \item \texttt{g\_cmdBacktraces}
    \item \texttt{extraWaitSemaphores} (guarded by \texttt{extraSemaphoreMutex})
    \item \texttt{semaphoresPendingDestroy}
    \item \texttt{deferredDestroys} (guarded by \texttt{deferredDestroyMutex})
    \item \texttt{buffersPendingAutoDestroy} (guarded by \texttt{buffersPendingMutex})
  \end{itemize}

  Each is paired with its own mutex, but lock ordering between the mutexes is
  documented only in comments, not enforced by the type system. Violations
  cause non-deterministic deadlocks. This also prevents unit testing and
  disallows multiple \texttt{VulkanApp} instances.

  \textbf{Severity Justification:} Violates AGENTS.md ``avoid global mutable
  state.'' Deadlocks are silent, non-deterministic, and notoriously hard to
  reproduce.

  \textbf{Recommendation:}
  Migrate all global state into \texttt{VulkanApp} member variables. Replace
  per-object mutexes with a single ordered lock or a lock-free structure.
  Alternatively, introduce a \texttt{SubmissionTracker} class and pass it
  by reference.

\item[\textbf{[G2]}] \textbf{\texttt{\_exit(1)} in Validation Callback}
  \hfill \severitybox{criticalcolor}{CRITICAL}

  \begin{tabular}{ll}
    \textbf{File:} & \texttt{vulkan/VulkanApp.cpp} \\
    \textbf{Line:} & 889 \\
    \textbf{Category:} & Robustness / Shutdown \\
  \end{tabular}

  \textbf{Description:}
  The \texttt{debugCallback} calls \texttt{\_exit(1)} on any warning or
  error-level validation message. \texttt{\_exit} bypasses all C++ destructors,
  leaking the GLFW window, Vulkan device, instance, and all GPU allocations.
  Even benign warnings during development — such as
  ``UNASSIGNED-BestPractices'' (which is explicitly silenced at line 864) —
  could theoretically trigger the hard exit if the suppression logic fails.

  \begin{lstlisting}
  // Line 887-889
  if (isWarningOrError) {
      _exit(1);
  }
  \end{lstlisting}

  \textbf{Severity Justification:} Prevents graceful GPU shutdown, makes
  debugging device-lost errors harder, and violates the AGENTS.md preference
  for ``descriptive log messages'' over hard exits.

  \textbf{Recommendation:}
  Replace \texttt{\_exit(1)} with a debug breakpoint (\texttt{\_\_builtin\_trap()}
  or \texttt{abort()}) so a core dump is generated but registered atexit
  handlers still run. Alternatively, log and continue; let the caller decide
  whether to exit.

\end{enumerate}

\newpage

\subsection{High Severity}

\begin{enumerate}[leftmargin=*]

\item[\textbf{[B7]}] \textbf{\texttt{ALL\_COMMANDS\_BIT} Used as Signal Stage Mask}
  \hfill \severitybox{highcolor}{HIGH}

  \begin{tabular}{ll}
    \textbf{File:} & \texttt{vulkan/VulkanApp.cpp} \\
    \textbf{Line:} & 2181, 2327, 4902, 4907 \\
    \textbf{Category:} & Vulkan API / Spec Violation \\
  \end{tabular}

  \textbf{Description:}
  Four semaphore signal operations use
  \texttt{VK\_PIPELINE\_STAGE\_2\_ALL\_COMMANDS\_BIT} as the \texttt{stageMask}
  in \texttt{VkSemaphoreSubmitInfo}:

\begin{lstlisting}
// Line 2181
signalSemaphoreInfo.stageMask = VK_PIPELINE_STAGE_2_ALL_COMMANDS_BIT;
\end{lstlisting}

  Per the Vulkan specification, \texttt{ALL\_COMMANDS\_BIT} is a \textit{wait}
  mask and is not valid in a semaphore signal operation. Although many
  implementations silently treat it as ``all stages have completed,'' the
  specification reserves the right to ignore it, and synchronization validation
  layers flag it as an error. This is particularly problematic at lines 4902
  and 4907 where the swapchain-present + frame-timeline semaphores are
  signaled — using an invalid stage mask there could interfere with WSI
  timing.

  \textbf{Recommendation:}
  For graphics submissions, use the union of all stages that produced output:
  \texttt{ALL\_GRAPHICS\_BIT} or the specific stage(s) (e.g.,
  \texttt{COLOR\_ATTACHMENT\_OUTPUT\_BIT | EARLY\_FRAGMENT\_TESTS\_BIT}).
  For compute submissions, use \texttt{COMPUTE\_SHADER\_BIT}.

\item[\textbf{[B8]}] \textbf{WaterRenderer Hardcodes 3 Instead of FRAMES Constant}
  \hfill \severitybox{highcolor}{HIGH}

  \begin{tabular}{ll}
    \textbf{File:} & \texttt{vulkan/renderer/WaterRenderer.cpp} \\
    \textbf{Line:} & 101, 105, 111, 140, 203, 207 \\
    \textbf{Category:} & Bug / Brittle Constant \\
  \end{tabular}

  \textbf{Description:}
  The header declares:
\begin{lstlisting}
// WaterRenderer.hpp:167
static constexpr uint32_t FRAMES = VulkanApp::MAX_FRAMES_IN_FLIGHT;
std::array<VkImage, FRAMES> waterDepthImages = {};
\end{lstlisting}

  However, every loop in the \texttt{.cpp} iterates with \texttt{frameIdx < 3}
  instead of \texttt{frameIdx < FRAMES}:
\begin{lstlisting}
// WaterRenderer.cpp:111
for (int frameIdx = 0; frameIdx < 3; ++frameIdx) { ... }
\end{lstlisting}

  If \texttt{MAX\_FRAMES\_IN\_FLIGHT} is ever changed from 3, these loops will
  silently access out-of-bounds memory via \texttt{std::array::operator[]}
  (no bounds check in release builds), corrupting adjacent member variables
  or triggering undefined behavior.

  \textbf{Recommendation:}
  Replace all \texttt{3} literals in loop bounds with \texttt{FRAMES}.

\item[\textbf{[B9]}] \textbf{Per-Cycle Backtrace Capture in Release Builds}
  \hfill \severitybox{highcolor}{HIGH}

  \begin{tabular}{ll}
    \textbf{File:} & \texttt{vulkan/VulkanApp.cpp} \\
    \textbf{Line:} & 2651--2665 \\
    \textbf{Category:} & Performance / CPU \\
  \end{tabular}

  \textbf{Description:}
  Every call to \texttt{allocatePrimaryCommandBuffer} for ring-slot command
  buffers captures a full stack backtrace via \texttt{backtrace()} and
  \texttt{backtrace\_symbols()}:

\begin{lstlisting}
void* bt_buf[32];
int bt_n = backtrace(bt_buf, 32);
char** bt_syms = backtrace_symbols(bt_buf, bt_n);
// ... format into std::string ...
g_cmdBacktraces[cmd] = bt_str;
\end{lstlisting}

  Since ring-slot command buffers are recycled every frame (up to 3 graphics
  + several async), this calls \texttt{malloc} inside
  \texttt{backtrace\_symbols} multiple times per frame even in release builds.
  Each call walks the ELF symbol table and formats strings — measurable CPU
  overhead on a hot path. The backtrace is only consumed on
  \texttt{VK\_ERROR\_DEVICE\_LOST}, which is rare in production.

  \textbf{Recommendation:}
  Guard with \texttt{\#ifndef NDEBUG} so the backtrace is only captured in
  debug builds. Alternatively, defer capture to the error path via
  \texttt{std::call\_once}.

\item[\textbf{[B10]}] \textbf{\texttt{signalRingSlotFence} Called Before Submit Result Check}
  \hfill \severitybox{highcolor}{HIGH}

  \begin{tabular}{ll}
    \textbf{File:} & \texttt{vulkan/VulkanApp.cpp} \\
    \textbf{Line:} & 2241--2244 \\
    \textbf{Category:} & Bug / Race Hazard \\
  \end{tabular}

  \textbf{Description:}
  In \texttt{submitCommandBufferAsync}, the ring-slot fence is signaled
  \textit{before} the \texttt{vkQueueSubmit2} result is checked for failure:

\begin{lstlisting}
VkResult submitRes = vkQueueSubmit2(graphicsQueue, 1, &submitInfo, fence);
if (submitRes == VK_SUCCESS) {
    signalRingSlotFence(commandBuffer, graphicsQueue);  // line 2243
}
if (submitRes != VK_SUCCESS) {
    // Cleanup path — but ring slot fence was already signaled
\end{lstlisting}

  If \texttt{vkQueueSubmit2} fails (with an error other than
  \texttt{DEVICE\_LOST}), the ring slot fence has already been signaled by
  \texttt{signalRingSlotFence}'s empty submit, making the ring slot appear
  available before the failed submission's command buffer has actually
  executed (it never will). This can cause the ring slot to be reused for a
  new recording that overlaps with the stale command buffer contents still
  on the GPU timeline.

  \textbf{Recommendation:}
  Move \texttt{signalRingSlotFence} after the success check, or guard it
  so it only fires on success.

\item[\textbf{[P6]}] \textbf{Processor Thread-Pool Tasks May Outlive Scene Data}
  \hfill \severitybox{highcolor}{HIGH}

  \begin{tabular}{ll}
    \textbf{File:} & \texttt{space/Processor.cpp} \\
    \textbf{Line:} & 40--70 (estimated) \\
    \textbf{Category:} & Memory Safety / Use-After-Free \\
  \end{tabular}

  \textbf{Description:}
  \texttt{Processor} enqueues tasks into the \texttt{ThreadPool} that capture
  raw \texttt{OctreeNode*} pointers. If the \texttt{Octree} or
  \texttt{ChunkManager} is destroyed while tasks are still pending or in
  flight, the captured pointers become dangling. The \texttt{ThreadPool}
  destructor joins all worker threads, but the octree may be destroyed
  independently (e.g., scene reload) while the pool is still running.

  \textbf{Recommendation:}
  Use \texttt{std::shared\_ptr} for octree nodes, or add a generation counter
  to each node that tasks check before dereferencing. Alternatively, drain
  the task queue before destroying the octree.

\item[\textbf{[P7]}] \textbf{\texttt{OctreeNode::getBlock} Returns Dangling Raw Pointer}
  \hfill \severitybox{highcolor}{HIGH}

  \begin{tabular}{ll}
    \textbf{File:} & \texttt{space/OctreeNode.cpp} \\
    \textbf{Category:} & Memory Safety \\
  \end{tabular}

  \textbf{Description:}
  \texttt{getBlock()} returns a raw \texttt{ChildBlock*} from a custom
  block allocator. When the allocator is reset (on full octree
  reinitialization), all previously returned pointers become dangling. The
  allocator provides no notification mechanism, so any holder of a
  \texttt{ChildBlock*} — such as \texttt{Processor} or
  \texttt{OctreeVisibilityChecker} — silently references freed memory.

  \textbf{Recommendation:}
  Use \texttt{std::shared\_ptr<ChildBlock>} with a custom deleter, or add an
  epoch/generation counter that is validated on every access.

\end{enumerate}

\newpage

\subsection{Medium Severity}

\begin{enumerate}[leftmargin=*]

\item[\textbf{[B4]}] \textbf{\texttt{vkAcquireNextImageKHR} Uses \texttt{UINT64\_MAX} Timeout}
  \hfill \severitybox{mediumcolor}{MEDIUM}

  \begin{tabular}{ll}
    \textbf{File:} & \texttt{vulkan/VulkanApp.cpp} \\
    \textbf{Line:} & 4767 \\
    \textbf{Category:} & Bug / Robustness \\
  \end{tabular}

  \textbf{Description:}
  \texttt{vkAcquireNextImageKHR} is called with \texttt{UINT64\_MAX} timeout
  (\textit{carried forward from report \#06}). If the window is minimized,
  the compositor is stalled, or the driver hangs, this blocks the main thread
  indefinitely. The swapchain is recreated on \texttt{VK\_ERROR\_OUT\_OF\_DATE\_KHR}
  but no timeout is enforced.

  \textbf{Recommendation:}
  Reduce the timeout to 1~s and handle \texttt{VK\_TIMEOUT} by yielding and
  retrying. Throttle to 1~FPS when minimized.

\item[\textbf{[S4]}] \textbf{Swapchain Barrier — \texttt{TOP\_OF\_PIPE} with \texttt{srcAccessMask = 0}}
  \hfill \severitybox{mediumcolor}{MEDIUM}

  \begin{tabular}{ll}
    \textbf{File:} & \texttt{vulkan/VulkanApp.cpp} \\
    \textbf{Line:} & 4965--4968 \\
    \textbf{Category:} & Synchronization \\
  \end{tabular}

  \textbf{Description:}
  The swapchain image transition (UNDEFINED → COLOR\_ATTACHMENT\_OPTIMAL) uses
  \texttt{srcStageMask = VK\_PIPELINE\_STAGE\_2\_TOP\_OF\_PIPE\_BIT} with
  \texttt{srcAccessMask = 0}. While the barrier is correct for UNDEFINED
  images (nothing to make visible), it is semantically misleading: the
  acquire semaphore already provides the execution dependency, so the entire
  barrier could be omitted. The mismatch between
  \texttt{TOP\_OF\_PIPE\_BIT} (earliest possible stage) and \texttt{0}
  (no access) confuses Vulkan validation tools and reviewers.

  \textbf{Recommendation:}
  Either remove the barrier entirely (the acquire semaphore suffices) or
  keep it and add a comment explaining it is vestigial.

\item[\textbf{[P8]}] \textbf{ThreadPool Stores Tasks as \texttt{std::function<void()>}}
  \hfill \severitybox{mediumcolor}{MEDIUM}

  \begin{tabular}{ll}
    \textbf{File:} & \texttt{space/ThreadPool.tpp} \\
    \textbf{Line:} & 25--45 (estimated) \\
    \textbf{Category:} & Performance / CPU \\
  \end{tabular}

  \textbf{Description:}
  Each enqueued task is type-erased into a \texttt{std::function<void()>},
  which heap-allocates for non-trivial captures. During terrain meshing,
  thousands of small tasks are dispatched per frame. The alloc/free churn
  adds measurable CPU overhead and fragments the heap.

  \textbf{Recommendation:}
  Use a small-buffer-optimized function type (e.g., \texttt{fx::function}
  with 32~B inline storage) or a custom task struct with an embedded
  virtual dispatch. For pointer-free captures, a template parameter pack
  + \texttt{std::invoke} avoids allocation entirely.

\item[\textbf{[P9]}] \textbf{StagingRingBuffer::allocate — Linear Scan Under Mutex}
  \hfill \severitybox{mediumcolor}{MEDIUM}

  \begin{tabular}{ll}
    \textbf{File:} & \texttt{vulkan/StagingRingBuffer.cpp} \\
    \textbf{Line:} & 47--80 \\
    \textbf{Category:} & Performance / Contention \\
  \end{tabular}

  \textbf{Description:}
  \texttt{allocate()} holds \texttt{mutex} while scanning the free list
  (\texttt{std::list}) linearly for a contiguous chunk. For a fragmented
  ring with many small allocations, this scan can take tens of microseconds
  while stalling other threads (e.g., async uploads). The function also
  blocks on a \texttt{unique\_lock} without timeout, so a stuck allocator
  blocks all staging.

  \textbf{Recommendation:}
  Replace the \texttt{std::list} free list with a free-block tree
  (e.g., \texttt{std::multiset} keyed by size) for O(log~n) allocation,
  or use a buddy allocator. Add a timeout to the \texttt{unique\_lock} and
  fall back to a fresh staging buffer if contention is high.

\item[\textbf{[P10]}] \textbf{\texttt{allocBuffer} / \texttt{bulkAllocBuffer} Lock Contention}
  \hfill \severitybox{mediumcolor}{MEDIUM}

  \begin{tabular}{ll}
    \textbf{File:} & \texttt{vulkan/renderer/IndirectRenderer.cpp} \\
    \textbf{Line:} & 207--248 \\
    \textbf{Category:} & Performance / Contention \\
  \end{tabular}

  \textbf{Description:}
  \texttt{bulkAlloc} holds \texttt{mutex} while inserting into
  \texttt{mergedVertices} and \texttt{mergedIndices}, which can be large
  (thousands of vertices). This blocks other operations — including
  \texttt{removeMesh}, \texttt{getActiveMeshInfos}, and
  \texttt{getMeshInfo} — from making progress. During level loading or
  streaming, this creates frame-time spikes.

  \textbf{Recommendation:}
  Use a reader-writer lock (\texttt{std::shared\_mutex}) where
  \texttt{getActiveMeshInfos} acquires shared access and
  \texttt{bulkAlloc}/\texttt{removeMesh} acquire exclusive access.
  Alternatively, batch inserts under the lock using \texttt{reserve} to
  minimize time under contention.

\item[\textbf{[B11]}] \textbf{Pipeline Cache Not Atomically Written}
  \hfill \severitybox{mediumcolor}{MEDIUM}

  \begin{tabular}{ll}
    \textbf{File:} & \texttt{vulkan/VulkanApp.cpp} \\
    \textbf{Line:} & 262--290 \\
    \textbf{Category:} & Bug / Data Integrity \\
  \end{tabular}

  \textbf{Description:}
  \texttt{savePipelineCache()} writes directly to
  \texttt{pipeline\_cache.bin}. If the application crashes mid-write, the
  file is truncated or corrupted, causing the next run to silently fall back
  to an empty cache. All pipelines must then be recompiled at load time, adding
  seconds to startup.

  \textbf{Recommendation:}
  Write to a temporary file (\texttt{pipeline\_cache.tmp}), then atomically
  rename over the target. This guarantees the cache is either fully written
  or the old version remains intact.

\item[\textbf{[G3]}] \textbf{\texttt{VulkanApp::run()} Swallows All Exceptions}
  \hfill \severitybox{mediumcolor}{MEDIUM}

  \begin{tabular}{ll}
    \textbf{File:} & \texttt{vulkan/VulkanApp.cpp} \\
    \textbf{Line:} & 5817--5821 \\
    \textbf{Category:} & Error Handling \\
  \end{tabular}

  \textbf{Description:}
  \texttt{mainLoop()} exceptions are caught by \texttt{catch (...)} with a
  generic message and then ignored:

\begin{lstlisting}
try {
    mainLoop();
} catch (...) {
    fprintf(stderr, "[VulkanApp] mainLoop threw, proceeding with cleanup\n");
}
cleanup();
\end{lstlisting}

  The exception type and \texttt{what()} string are lost. For critical
  failures during rendering (descriptor pool exhaustion, device loss), the
  engine proceeds to \texttt{cleanup()} with the device in an unknown state,
  potentially causing secondary crashes during teardown.

  \textbf{Recommendation:}
  Catch \texttt{std::exception} and log \texttt{e.what()} before cleanup.
  Keep a bare \texttt{catch (...)} as a last resort but still log the
  address of the exception record.

\item[\textbf{[P11]}] \textbf{Broken Double-Checked Locking on Vegetation Pipeline}
  \hfill \severitybox{mediumcolor}{MEDIUM}

  \begin{tabular}{ll}
    \textbf{File:} & \texttt{vulkan/VulkanApp.cpp} \\
    \textbf{Line:} & 675--815 \\
    \textbf{Category:} & Concurrency / Waste \\
  \end{tabular}

  \textbf{Description:}
  The \texttt{ensureVegetationComputePipeline()} function locks the mutex
  for the initial check, then releases it before the expensive pipeline
  creation work:

\begin{lstlisting}
// Line 677-680 — lock + initial check
{
    std::lock_guard<std::mutex> lk(vegComputeMutex);
    if (vegComputePipeline != VK_NULL_HANDLE) return true;
}
// Line 682-774 — EXPENSIVE init, OUTSIDE lock
// ...
// Line 777-813 — re-acquire lock for commit
{
    std::lock_guard<std::mutex> lk(vegComputeMutex);
    if (vegComputePipeline == VK_NULL_HANDLE) {
        // commit
    } else {
        // clean up duplicates
    }
}
\end{lstlisting}

  This allows a second thread to also pass the initial check and begin a
  duplicate initialization before the first thread commits. The duplicate
  work is then discarded (lines 806--813). The proper pattern would be to
  check outside the lock for the fast path:
\begin{lstlisting}
if (vegComputePipeline != VK_NULL_HANDLE) return true;  // fast path, no lock
{
    std::lock_guard... lk(vegComputeMutex);
    if (vegComputePipeline != VK_NULL_HANDLE) return true;  // double-check
    // ... initialize under lock ...
}
\end{lstlisting}

  \textbf{Recommendation:}
  Restructure to the proper double-checked locking pattern so the fast path
  avoids the mutex entirely and the initialization is serialized.

\end{enumerate}

\newpage

\subsection{Low Severity}

\begin{enumerate}[leftmargin=*]

\item[\textbf{[P5]}] \textbf{Per-Frame Descriptor Writes in SceneRenderer}
  \hfill \severitybox{lowcolor}{LOW}

  \begin{tabular}{ll}
    \textbf{File:} & \texttt{vulkan/renderer/SceneRenderer.cpp} \\
    \textbf{Line:} & 764--773 \\
    \textbf{Category:} & Performance / CPU \\
  \end{tabular}

  \textbf{Description:}
  Descriptor sets for uniform buffers and image samplers are updated every
  frame via \texttt{vkUpdateDescriptorSets} (carried forward from report \#06).
  For resources that change infrequently (static textures, constant buffers),
  this generates unnecessary CPU work and driver validation overhead.

  \textbf{Recommendation:}
  Separate descriptor sets into static (updated once) and dynamic (updated
  per frame). Use the new \texttt{DescriptorWriter} builder for the dynamic
  portion.

\item[\textbf{[C2]}] \textbf{Binary Semaphores Used Alongside Timeline}
  \hfill \severitybox{lowcolor}{LOW}

  \begin{tabular}{ll}
    \textbf{File:} & \texttt{vulkan/VulkanApp.cpp} \\
    \textbf{Line:} & 4799 \\
    \textbf{Category:} & Compliance \\
  \end{tabular}

  \textbf{Description:}
  The present submit signals both a binary semaphore (required by WSI) and
  the frame timeline semaphore (carried forward from report \#06). This dual
  signal path complicates synchronization reasoning but is inherently
  required by the Vulkan WSI extension.

  \textbf{Recommendation:}
  Acceptable as-is. Add a comment at the submission site explaining why both
  are present.

\item[\textbf{[C3]}] \textbf{SPIR-V Targets \texttt{vulkan1.1} Instead of 1.6}
  \hfill \severitybox{lowcolor}{LOW}

  \begin{tabular}{ll}
    \textbf{File:} & \texttt{Makefile} \\
    \textbf{Category:} & Compliance \\
  \end{tabular}

  \textbf{Description:}
  Shaders target \texttt{vulkan1.1} (SPIR-V 1.3). The engine uses Vulkan 1.3
  features (carried forward from report \#06). Bumping to \texttt{vulkan1.3}
  (SPIR-V 1.6) enables 8/16-bit types, scalar block layout, and 64-bit
  integers in shader code.

  \textbf{Recommendation:}
  Change \texttt{--target-env vulkan1.1} to \texttt{vulkan1.3} (or
  \texttt{vulkan1.4} when available) in the \texttt{glslc} flags.

\item[\textbf{[B6]}] \textbf{Resource Manager Tracks \texttt{VkRenderPass} Unnecessarily}
  \hfill \severitybox{lowcolor}{LOW}

  \begin{tabular}{ll}
    \textbf{File:} & \texttt{vulkan/VulkanResourceManager.cpp} \\
    \textbf{Line:} & 176, 275, 293, 488 \\
    \textbf{Category:} & Dead Code \\
  \end{tabular}

  \textbf{Description:}
  \texttt{VulkanResourceManager} tracks \texttt{VkRenderPass} handles, but
  no render passes are ever created (carried forward from report \#06).

  \textbf{Recommendation:}
  Remove the tracking code.

\item[\textbf{[G4]}] \textbf{Duplicate \texttt{return VK\_FALSE} in \texttt{debugCallback}}
  \hfill \severitybox{lowcolor}{LOW}

  \begin{tabular}{ll}
    \textbf{File:} & \texttt{vulkan/VulkanApp.cpp} \\
    \textbf{Line:} & 891--892 \\
    \textbf{Category:} & Code Quality \\
  \end{tabular}

  \textbf{Description:}
  Two consecutive \texttt{return VK\_FALSE;} statements. The second is dead
  code.

\begin{lstlisting}
    return VK_FALSE;
    return VK_FALSE;  // dead
}
\end{lstlisting}

  \textbf{Recommendation:}
  Remove the second \texttt{return}.

\item[\textbf{[P12]}] \textbf{\texttt{getActiveMeshInfos()} Copies Vector Every Call}
  \hfill \severitybox{lowcolor}{LOW}

  \begin{tabular}{ll}
    \textbf{File:} & \texttt{vulkan/renderer/IndirectRenderer.cpp} \\
    \textbf{Line:} & 1314--1321 \\
    \textbf{Category:} & Performance / Allocations \\
  \end{tabular}

  \textbf{Description:}
  Returns a \texttt{std::vector<MeshInfo>} by value, allocating and copying
  every active mesh entry. Only used from the debug UI, so the impact is
  limited, but during frame spikes (thousands of debug toggles) it adds
  allocator pressure.

  \textbf{Recommendation:}
  Accept an output iterator or return a vector of indices to the internal
  storage.

\end{enumerate}

\newpage

\section{Comparison with Previous Report}

\begin{center}
\begin{tabular}{lccccc}
\toprule
\textbf{Category} & \textbf{Report \#06} & \textbf{Resolved} & \textbf{Carried} & \textbf{New} & \textbf{Total Open} \\
\midrule
Bugs (B)           & 6  & 4 & 2  & 4 & 6  \\
Performance (P)    & 5  & 3 & 2  & 4 & 6  \\
Compliance (C)     & 3  & 1 & 2  & 0 & 2  \\
Synchronization (S)& 4  & 3 & 1  & 0 & 1  \\
General (G)        & —  & — & —  & 4 & 4  \\
\midrule
\textbf{Total}     & \textbf{18} & \textbf{11} & \textbf{7} & \textbf{12} & \textbf{19} \\
\bottomrule
\end{tabular}
\end{center}

\textit{Note:} The 11 resolved issues were all from the bug, performance,
compliance, and synchronization categories. The 7 carried-forward issues are
lower-severity items that were not prioritized. The 12 new findings are
concentrated in concurrency safety (\texttt{space/}) and global architecture
(\texttt{VulkanApp.cpp}).

\newpage

\section{Overall Classification}

The engine has made substantial progress since report \#06. The 11 resolved
issues include all previously flagged critical bugs (async ring wrap, device-lost
cleanup, buffer leaks) and several high-impact performance antipatterns
(\texttt{deviceWaitIdle}, \texttt{ALL\_COMMANDS\_BIT}, \texttt{vkAllocateMemory}).
The introduction of \texttt{DescriptorWriter}, \texttt{DescriptorAllocator},
\texttt{GraphicsPipelineConfig}, \texttt{TrackedHandle<T>}, and
\texttt{RendererUtils::transitionImageLayout} has measurably improved code
quality and reduced boilerplate by over 600 lines.

However, three systemic concerns remain:

\begin{enumerate}
  \item \textbf{Global mutable state (G1) is the single highest-risk item.}
        The eight file-scope data structures with independent mutexes create
        a fragile locking graph that can deadlock non-deterministically.
        This should be the top refactoring priority.
  \item \textbf{The \texttt{\_exit(1)} in the validation callback (G2)}
        defeats all RAII guarantees and makes every validation warning
        potentially fatal during development. Replacing it with a
        breakpoint or abort would preserve debuggability without
        sacrificing cleanup.
  \item \textbf{Lifetime management in \texttt{space/} (P6, P7)} exposes the
        engine to use-after-free bugs during scene reload. The raw-pointer
        patterns in the octree allocator and thread-pool task dispatch need
        modernization to \texttt{shared\_ptr} or generation counters.
\end{enumerate}

\subsection*{Overall Severity}

\medskip
\noindent\severitybox{highcolor}{ADEQUATE}

\begin{quote}
  The engine is significantly more robust than at report \#06. The core
  rendering path (dynamic rendering, Synchronization2, timeline semaphores,
  VMA suballocation) follows modern best practices and is validation-clean
  in common paths. \textbf{2 critical, 6 high, and 8 medium issues remain},
  concentrated in the global state layer, validation handler, and space
  partitioning lifetime management. None of these are expected to manifest
  in normal operation, but they represent real hazards under edge cases:
  concurrent scene reload, rapid device-lost recovery, or
  \texttt{MAX\_FRAMES\_IN\_FLIGHT} adjustments.
\end{quote}

\subsection*{Remediation Priority}

\begin{enumerate}[noitemsep]
  \item[\textbf{P0}] \textbf{This sprint:} G1 (global mutable state refactor),
        G2 (\texttt{\_exit(1)} → breakpoint), B8 (WaterRenderer 3→FRAMES).
  \item[\textbf{P1}] \textbf{Next sprint:} B7 (ALL\_COMMANDS\_BIT signal mask),
        B10 (ring-slot fence ordering), P6/P7 (octree lifetime).
  \item[\textbf{P2}] \textbf{This month:} B4 (acquire timeout), B9 (backtrace
        guard), B11 (atomic cache write), P8-P10 (lock contention).
  \item[\textbf{P3}] \textbf{Backlog:} P5 (descriptor writes), C2/C3 (compliance),
        B6 (dead code), G4 (duplicate return), P12 (copy elision).
\end{enumerate}

\newpage

\section{Fix Instructions Summary}

\begin{enumerate}[noitemsep]
  \item G1 — Move 8 global structures into \texttt{VulkanApp} members;
        unify mutex ordering with a single ordered lock or lock-free
        alternative.
  \item G2 — Replace \texttt{\_exit(1)} with \texttt{\_\_builtin\_trap()} or
        \texttt{abort()}.
  \item G3 — Add \texttt{catch (std::exception\& e)} before
        \texttt{catch (...)} and log \texttt{e.what()}.
  \item G4 — Remove duplicate \texttt{return VK\_FALSE;} at line 892.
  \item B4 — Reduce \texttt{vkAcquireNextImageKHR} timeout to 1~s; handle
        \texttt{VK\_TIMEOUT} by yielding.
  \item B6 — Remove \texttt{VkRenderPass} tracking from
        \texttt{VulkanResourceManager}.
  \item B7 — Replace \texttt{ALL\_COMMANDS\_BIT} in signal stage masks with
        the narrowest appropriate stage (\texttt{ALL\_GRAPHICS\_BIT},
        \texttt{COMPUTE\_SHADER\_BIT}, etc.).
  \item B8 — Replace \texttt{3} with \texttt{FRAMES} (i.e.,
        \texttt{VulkanApp::MAX\_FRAMES\_IN\_FLIGHT}) in all
        \texttt{WaterRenderer.cpp} loop bounds.
  \item B9 — Guard \texttt{backtrace()}/\texttt{backtrace\_symbols()} with
        \texttt{\#ifndef NDEBUG}.
  \item B10 — Move \texttt{signalRingSlotFence} after the success check.
  \item B11 — Write pipeline cache to a temp file, then atomically rename.
  \item P5 — Use static descriptor sets for immutable resources.
  \item P6 — Use \texttt{shared\_ptr} or generation counters in octree
        task dispatch.
  \item P7 — Add epoch validation to \texttt{ChildBlock*} accessors.
  \item P8 — Use small-buffer-optimized function type in ThreadPool.
  \item P9 — Replace free-list scan with size-indexed tree or buddy allocator.
  \item P10 — Use \texttt{std::shared\_mutex} for IndirectRenderer mesh
        access.
  \item P11 — Restructure to proper double-checked locking (check outside lock).
  \item P12 — Return by span or iterator pair instead of vector copy.
  \item C2 — Add explanatory comment at the dual-signal submit site.
  \item C3 — Bump SPIR-V target to \texttt{vulkan1.3+} in \texttt{glslc} flags.
  \item S4 — Remove vestigial swapchain barrier or add explanatory comment.
\end{enumerate}

\end{document}
